\section{From false acceptance to theft}
\label{sec:theft-state}

A false-acceptance theorem is necessary for theft resistance, but it is not
the whole claim.  Theft also depends on how a victim note is chosen, what the
attacker may observe before acting, how note and nullifier hashes are used,
and whether the chain records a successful spend exactly once.  The Lean
development treats these questions separately.

\subsection{A fixed victim}

The simplest game fixes one live note, its depth-20 position, and its Merkle
siblings before the attacker acts.  The target contains a valid opening
\((k_\nu,r_{\mathrm{in}},v_{\mathrm{in}},a)\).  Its public note commitment,
anchor, and nullifier are computed with the deployed owner, note, nullifier,
and node functions.

The event \code{FixedVictimAttackEvent} covers two ways of targeting that
note.  An accepted proof either uses the victim's nullifier, or it presents a
leaf at the victim's exact position under the victim anchor.  This definition
also includes the honest owner.  In a security game, an honest opening is
therefore classified as credential recovery rather than silently excluded.

Lean proves the following complete case split.

\begin{theorem}[Fixed-victim classification]
\label{thm:fixed-victim}
Every accepted fixed-victim attack implies at least one of:
\begin{enumerate}[leftmargin=*]
  \item extraction failed to produce a witness satisfying the spend relation;
  \item the extracted witness contains the victim's secret and nullifier
        randomness;
  \item a different secret/randomness pair has the victim's nullifier;
  \item a different valid note opening has the victim's input commitment; or
  \item a different leaf at the victim's exact tree position reaches the
        victim anchor.
\end{enumerate}
In the last case, the two paths expose distinct ordered node inputs with the
same node-hash output at some level.
\end{theorem}

The pointwise result is
\code{fixed_victim_attack_implies_mathematical_failure}.  Its set form is
\code{fixed_victim_attack_subset_mathematical_failures}, and
\code{fixed_victim_attack_measure_le_five_failures} gives the corresponding
union bound.  None of these theorems assigns a numerical probability to key
recovery or to a Poseidon2 collision or second preimage.

The position condition matters.  A different direction-bit sequence denotes
a different Merkle position.  Equality of two roots reached at different
positions is not, by itself, a collision witness for one node call.  The
reduction therefore requires the same position before it derives a concrete
node collision.  This avoids an invalid global injectivity assumption for a
compressing hash.

\subsection{First fraudulent spend after public observations}

A fixed victim does not cover an attacker who first observes other proofs.
The module \code{V5AdaptiveObservedTheftGame} models an arbitrary finite,
ordered history of public proof records.  The attacker sees the public
projection of that history, chooses a spend attempt, and may depend on the
entire observed prefix.  The success event is the first fraudulent spend
selected in this experiment.

The theorem
\code{adaptive_first_fraudulent_spend_implies_mathematical_failure} repeats
the five-way classification with a history-dependent extractor.  Theorems
\code{adaptive_first_fraudulent_spend_subset_total_final_failure} and
\code{deployed_adaptive_first_fraudulent_spend_subset_final_or_setup} connect
that classification to the final failure ledger and retain a separate event
for an invalid or ambiguous victim setup.

This is a finite-history theorem, not an unlimited-attempt guarantee.  If a
caller permits several final attack attempts, it must bound their number or
prove an appropriate state-restoration statement.  It must also show that the
soundness and extraction premises hold for the distribution induced by the
observed history.

\subsection{Raw theft accounting}

For a submitted accepted proof, the attacker has already completed the work
search.  The raw accounting must therefore omit the 37-bit normalization used
in the budget theorem of \cref{thm:conditional-soundness}.  Lean proves the
conditional inequality
\begin{equation}
  \Prob[\text{first fraudulent spend}]
  \leq 2^{-75}+B_{\mathrm{external}}+
    \Prob[\text{invalid victim setup}]
  \label{eq:raw-theft-75}
\end{equation}
as
\code{deployed_adaptive_first_fraudulent_spend_probability_le_two_pow_neg_75}.
The external budget in this theorem is a premise; it includes production,
primitive, credential, and runtime terms rather than proving them small.

A later width-19 theorem substitutes checked arithmetic for the raw batching
term.  Its conclusion has the form
\begin{align}
  \Prob[\text{first fraudulent spend}] \leq {}& 2^{-70}
    + B_{\mathrm{nonterminal\ statement}}
    + \Prob[\text{hash/Merkle residual}] \notag\\
    &+ B_{\mathrm{crypto}}+B_{\mathrm{credential}}
    +B_{\mathrm{runtime}}+\Prob[\text{invalid victim setup}].
  \label{eq:raw-theft-70}
\end{align}
The exact theorem is
\code{deployed_qm31_adaptive_theft_probability_le_two_pow_neg_70_plus_remaining}.
It is an accounting result with visible remaining terms.  It is not a
numerical deployed theft bound.

\subsection{Spent-marker transition}

The source-shaped state model distinguishes account validation, proof
acceptance, a recheck of mutable pool state, spent-marker creation, and the
pool-root update.  A successful transition requires five accounts in the
specified roles, the expected program owners and writable flags, the expected
derived marker address, derivation bump 255, a usable marker account, and the
current pool anchor and sequence number.

The model represents two permitted pre-spend marker states.  A fresh
system-owned empty account may be prepared for the marker, and an already
prepared empty program-owned marker may be written.  A marker that already
contains a spent record is rejected.  On success, the transition records the
nullifier, writes the new pool anchor, and increments the sequence number.

The principal results are:
\begin{itemize}[leftmargin=*]
  \item \code{v5_source_success_implies_requirements}: every modeled success
        satisfies all account, address, bump, proof, marker, and mutable-state
        checks;
  \item \code{v5_source_success_exact_state}: the returned state is exactly
        the modeled marker and pool update;
  \item \code{v5_source_rejects_preexisting_marker}: an existing marker
        cannot be overwritten;
  \item \code{v5_source_success_exact_trace}: a successful trace ends with
        the mutable-state recheck, marker write, and pool write; and
  \item \code{no_sequential_success_for_the_same_nullifier}: two modeled
        sequential successes cannot use the same nullifier.
\end{itemize}

The address case is slightly stronger than a replay-only statement.  If two
different nullifiers derive to the same marker address, the first marker still
occupies the account and the second transition cannot overwrite it.  Such an
address alias remains a named failure in the full reduction because the
connection between production address derivation and the abstract model is a
separate assumption.

The marker prevents a second successful spend after the first has committed.
It cannot prevent the first successful spend from being fraudulent.  That is
why extraction, credentials, hash binding, and the first-fraudulent-spend
game remain necessary.

\subsection{Close, refund, rollback, and finality}

Proof-account closing is modeled as a later transaction.  The close path
checks the proof and refund account owners, signer and writable flags,
distinct addresses, an open proof account, and a positive balance.  On
success it invalidates the proof account, credits the entire modeled proof
balance to the supplied refund account, and leaves the pool and marker state
unchanged.  This is proved by
\code{close_success_refunds_exact_recipient_and_preserves_spend}.
The ordering theorem
\code{v5_then_close_trace_has_state_writes_before_refund} places the committed
marker and pool writes before the later proof invalidation and refund.

The equation \code{rejected_outcome_rolls_back_in_model} is true by the
definition of the abstract visible-state projection.  It does not prove
Solana rollback.  Likewise, the Lean transition does not prove account locks,
system-program behavior, or persistence at finalized commitment.  These are
fields of \code{ProductionRuntimePremises}; their negations appear as named
events in the final ledger.

At the source boundary, \code{ExactCurrentRustStateEquality} and
\code{ExactRustCloseEquality} remain the exact equalities needed to transfer
the source-shaped transition theorems to all successful production runs.
Recorded mainnet state supports these equalities for one execution, but one
execution is not their universal proof.

\begin{table}[t]
\centering
\small
\caption{Theft and state results.}
\label{tab:theft-state-anchors}
\begin{tabularx}{\textwidth}{@{}YY@{}}
\toprule
Purpose & Lean theorem or definition \\
\midrule
Five-way fixed-victim split &
\code{fixed_victim_attack_implies_mathematical_failure} \\
Eight-term chain-level union bound &
\code{deployed_attack_measure_le_eight_failures} \\
Observed-history classification &
\code{adaptive_first_fraudulent_spend_implies_mathematical_failure} \\
Raw 75-bit conditional accounting &
\code{deployed_adaptive_first_fraudulent_spend_probability_le_two_pow_neg_75} \\
Width-19 raw conditional accounting &
\code{deployed_qm31_adaptive_theft_probability_le_two_pow_neg_70_plus_remaining} \\
Exact modeled spend state & \code{v5_source_success_exact_state} \\
Existing marker rejection & \code{v5_source_rejects_preexisting_marker} \\
Close and refund behavior &
\code{close_success_refunds_exact_recipient_and_preserves_spend} \\
\bottomrule
\end{tabularx}
\end{table}
